% SPDX-License-Identifier: Apache-2.0
% covers: AxiWb
\datasheet{AxiWb}{A bridge from the peripheral end of an AXI4 link to the host lines of a pipelined Wishbone bus.}

\dsfacts{\code{lib/parts/src/bus/wb.rs}}{\code{txhdl\_parts::bus::wb::AxiWb}}%
{\code{//docs:axi}, \code{//docs:vreteno}}

\subsection*{Function}

A controller written elsewhere often speaks Wishbone rather than AXI.
\code{AxiWb} puts such a controller on a link. It is a peripheral
client written as hardware. On one side it holds the channel ends an
\code{AxiPer} gives a client: requests on \code{req} and write beats on
\code{wd} in, answers on \code{ans} and read beats on \code{rb} out.
On the other side it drives the host lines of a pipelined Wishbone
bus, \code{cyc}, \code{stb}, \code{we}, \code{adr}, \code{dat} and
\code{sel}, and reads \code{stall}, \code{ack} and the peripheral's
word, \code{rdat}.

It takes one single-beat burst at a time and puts it on the lines as
one request. The DDR3 controller of Vreteno's board is behind one.

\subsection*{Parameters}

\begin{center}\footnotesize
\begin{tabular}{@{}lp{0.7\textwidth}@{}}
\toprule
Parameter & Meaning \\
\midrule
\code{A} & The link's address width, in bits. \\
\code{I} & The link's identifier width, in bits. \\
\code{AW} & The width of the Wishbone word address, in bits. \\
\bottomrule
\end{tabular}
\end{center}

Words are 32 bits with four lanes, fixed in the source. The tables
below use \code{AxiWb<32, 2, 28>}. There \code{AW} is \code{ddr3::AW},
28 bits, the word address of the board's DDR3 memory.

\dstables{AxiWb}

\subsection*{Behaviour}

\code{stage} says where the request is. The source names five values.

\begin{itemize}
\item \textbf{0, waiting.} A request on \code{req} is taken when one
offers. The bridge keeps its identifier and its word address, and sets
\code{wsel} to all four lanes. A read goes to stage 2 at once, and a
write to stage 1.
\item \textbf{1, a write waiting for its beat.} The beat on \code{wd}
is taken when it offers. Its \code{data} goes to \code{wdat} and its
\code{strb} to \code{wsel}. Then stage 2.
\item \textbf{2, offered.} \code{cyc} and \code{stb} are high. A cycle
with \code{stall} low is the one the peripheral takes the request in,
and the stage moves to 3.
\item \textbf{3, taken.} \code{cyc} alone is high until \code{ack}.
An \code{ack} in the cycle of the take counts too. On \code{ack} the
stage moves to 4, and a read keeps the word on \code{rdat} in
\code{wdat}.
\item \textbf{4, answering.} A read sends one beat on \code{rb}: the
identifier, \code{wdat}, \code{Okay}, and \code{last} high. A write
sends an \code{Answer} of \code{Okay} on \code{ans}. Each waits for
room on its channel, then the stage returns to 0.
\item \code{we} is high unless the request is a read. \code{adr},
\code{dat} and \code{sel} are \code{wadr}, \code{wdat} and \code{wsel}.
Every line depends only on the bridge's registers, so nothing
combinational passes from the link to the lines or back.
\item The Wishbone address is a word address: the burst's byte address
shifted down by two and cut to \code{AW} bits. A peripheral decoded at
a region of the link's map sees offsets from the region's base, as
long as the base lies above those bits. A write at
\code{0x4000\_0010} lands in word 4.
\end{itemize}

\subsection*{Verification}

\begin{itemize}
\item Two tests in \code{bus::wb}, in \code{//lib/parts:parts\_test},
run the bridge on a link, behind \code{AxiHost} and \code{AxiPer},
against \code{WbMem}, a simulated Wishbone memory. The first,
\code{a\_read\_gets\_what\_a\_write\_put}, writes and reads one word.
The second, \code{slow\_answers\_and\_a\_calibration}, writes and
reads back 24 words at pseudorandom addresses, against a memory that
stalls for its first 40 cycles and answers after 5.
\item \code{ex\_axi\_wb} writes three words and reads them back, against
a memory that stalls for 10 cycles and answers after 2. The build
simulates the netlist of \code{AxiWb<32, 2, 28>} against that run under
nvc and Verilator: \code{//docs:axi\_wb\_sim\_axi\_wb\_tb\_test} and
\code{//docs:axi\_wb\_vsim\_test}.
\item In \code{//ddr3:ddr3\_test}, \code{words\_go\_in\_and\_come\_back}
runs \code{Ddr3Per}, this bridge with the controller's simulation
model behind it, and
\code{the\_controller\_is\_instantiated\_and\_not\_written} checks that
its netlist writes the bridge's module.
\item \code{//cpu/vreteno/board/sim:board\_test}, a manual target,
simulates the whole board under Vivado's simulator.
\end{itemize}

\subsection*{Used by}

\code{Ddr3Per} in \code{ddr3/src/lib.rs}, as its field \code{bridge},
and so the board's DDR3 memory at \code{0x4000\_0000}. The example
\code{ex\_axi\_wb}.

\subsection*{Limits}

It handles single-beat bursts only. It takes one write beat per
request, and answers a read with one beat marked \code{last}. It holds
one request at a time. It always answers \code{Okay}, and it reads no
error line from the peripheral. Words are 32 bits with four lanes.
Address bits above the \code{AW} bits it keeps are dropped.
